\section{Cross-Lens Study Protocol, Controls, and Per-Family Results}
\label{app:cross-lens}

This appendix documents the corpus conditionality study of
\S\ref{subsec:cross-lens}.

\paragraph{Lenses.}
We use the reference implementation of the Jacobian lens
\citep{gurnee2026workspace}, which reads an intermediate hidden state by
transporting it through an averaged Jacobian into the space consumed by
the unembedding and decoding through the LM head. On Qwen3.5-9B we fit
two lenses that differ only in fitting corpus, 100 seeded English
C4 prompts against 100 seeded Chinese C4 prompts, with identical seeds,
positions, and procedure. Transport is evaluated at layers
\(\{2,5,7,10,12,14,17,19,21,24,26,29\}\) of 31 at the final prompt
position. The cross-lens comparison uses layers 21, 24, 26, and 29,
where the two lenses' token outputs can diverge, while both lenses converge at
the final layer. Composition is validated numerically, since
the decoded states \(\rstate=T_\ell(h_\ell)\) reproduce the lens's own output scores
to within 0.2\%.

\paragraph{SRP basis and comparison operation.}
Because the implemented lens feeds the unembedding directly, the score
each lens reports is a selected readout score on \(\rstate\),
and SRP decomposes it in the shared basis fitted from the weights alone
(\S\ref{sec:method-comparing-accounts}). The dictionary is an
independently trained \(k{=}128\) (seed 0) readout SAE for
Qwen3.5-9B, a different operating point from the
\(32\times\)/\(k{=}256\) configuration of the main tables, disclosed here per the
rule of Appendix~\ref{app:basis-stability} for comparisons across recipes. No translation
pairs, language labels, or fitting corpora enter the basis, which is
fitted to the rows of the LM head only. For each prompt, each lens's
reported score is decomposed at layers 21, 24, 26, and 29. The dominant
feature at a layer is the one whose signed contribution is largest in
absolute value. We count a prompt as agreement when the same feature id
dominates under both lenses in at least two of the four layers.

\paragraph{Prompts.}
Eighty evaluation prompts span seven families, namely Chinese antonym (14),
English cloze (10), Chinese cloze (18), Chinese exemplar (8), Chinese
factual recall (10), and translation in both directions (10 each). A
further 12 control prompts (6 digit, 6 proper noun) have an expected
surface form that is language-invariant. Two null comparisons price in loose
matching. Under the unrelated token null, the dominant
feature of each reading is compared against the decomposition of an unrelated
token under the same lens on the same state (159 comparisons). Under the shuffled pairing
null, the two lenses' readings are re-paired across different prompts
(240 comparisons).

\begin{table}[!tbp]
\caption{\textbf{Cross-lens agreement on the dominant feature by prompt
family.} Agreement means the same SRP feature is dominant under the English and
the Chinese lens in at least two of layers 21, 24, 26, and 29, with 95\%
binomial CIs.
``Lens-only'' is the surface baseline, counting prompts where the English lens
reports top-1 tokens in Latin script and the Chinese lens in CJK script
(modal script across the same four layers), i.e.\ where token outputs
alone suggest disagreement.
Control families (below the rule) are excluded from the 77/80 headline.}
\label{tab:app-cross-lens-families}
\centering
\scriptsize
\setlength{\tabcolsep}{4.5pt}
\renewcommand{\arraystretch}{1.08}
\begin{tabularx}{\linewidth}{@{}>{\RaggedRight\arraybackslash}Xccc@{}}
\toprule
Family & Agreement & 95\% CI & Lens-only split \\
\midrule
Antonym (ZH) & 14/14 & [0.78, 1.00] & 4/14 \\
Cloze (EN) & 10/10 & [0.72, 1.00] & 7/10 \\
Cloze (ZH) & 16/18 & [0.67, 0.97] & 6/18 \\
Exemplar (ZH) & 8/8 & [0.68, 1.00] & 7/8 \\
Factual recall (ZH) & 9/10 & [0.60, 0.98] & 9/10 \\
Translation EN\(\to\)ZH & 10/10 & [0.72, 1.00] & 6/10 \\
Translation ZH\(\to\)EN & 10/10 & [0.72, 1.00] & 0/10 \\
\midrule
All cross-lens & 77/80 & [0.90, 0.99] & --- \\
\midrule
Digit controls & 5/6 & [0.44, 0.97] & 0/6 \\
Proper noun controls & 5/6 & [0.44, 0.97] & 0/6 \\
\bottomrule
\end{tabularx}
\end{table}

\paragraph{Results.}
The same SRP feature is dominant under both lenses on 77 of 80
cross-lens prompts (96\%, CI [0.90, 0.99]), uniformly across families
(\Cref{tab:app-cross-lens-families}). Within a single lens, one feature
carries both surface forms (e.g.\ the English and Chinese realizations
of the same concept) on 61/80 (English lens) and 62/80 (Chinese lens)
prompts. Both null floors are empty, at 0/159 matches under the unrelated
token null and 0/240 under the shuffled pairing null, so the headline agreement is
not an artifact of a permissive matching criterion. The ``Lens-only''
column shows the surface picture the decomposition corrects, since on
factual recall prompts the two lenses report different scripts on 9/10
prompts while the dominant feature agrees on 9/10. The worked example is
the prompt \citet{gurnee2026workspace} use for their multilingual
illustration, ``the opposite of
\begin{CJK*}{UTF8}{gbsn}小\end{CJK*}''. The lens fitted on English reports
\texttt{large}/\texttt{big} at layers 24, 26, and 29, while the lens
fitted on Chinese reports
\begin{CJK*}{UTF8}{gbsn}大的/大\end{CJK*}, and both readings are carried by
the same readout feature, whose top unembedding rows are the
\begin{CJK*}{UTF8}{gbsn}大\end{CJK*}~family.
\Cref{tab:app-cross-lens-antonym-layers} gives that case layer by layer.
The shared feature f112 is the
top contributor to \begin{CJK*}{UTF8}{gbsn}大\end{CJK*} under both
fitted lenses at each comparison layer (\(+12.5\) and \(+13.4\) at layer
24, \(+14.5\) and \(+17.6\) at layer 29) and is also the top
contributor to \texttt{big} under both, so the two surface forms are
carried by one coordinate within each lens as well as across the pair.

\begin{table}[!tbp]
\caption{\textbf{The antonym case, layer by layer.} Top-1
token and its softmax share under each lens on identical Qwen3.5-9B
hidden states \(h_\ell\) for the prompt of \citet{gurnee2026workspace}
(\begin{CJK*}{UTF8}{gbsn}``小''的反义词是\end{CJK*}, and the model answers
\begin{CJK*}{UTF8}{gbsn}大\end{CJK*}). The two Jacobian lenses differ
only in fitting corpus.}
\label{tab:app-cross-lens-antonym-layers}
\centering
\footnotesize
\setlength{\tabcolsep}{5pt}
\renewcommand{\arraystretch}{1.08}
\begin{CJK*}{UTF8}{gbsn}
\begin{tabular}{@{}lcc@{}}
\toprule
Layer & EN-fitted & ZH-fitted \\
\midrule
24 & \texttt{large} 40\% & 大的 36\% \\
26 & \texttt{large} 26\% & 大的 73\% \\
29 & \texttt{big} 58\% & 大 93\% \\
final & 大 98\% & 大 98\% \\
\bottomrule
\end{tabular}
\end{CJK*}
\end{table}

\paragraph{A second language pair.}
The English--German bank removes the script confound, since both languages use
the Latin script, so agreement cannot come from the script alone. Ninety
cross-lens prompts in seven families mirror the English--Chinese bank
(German antonym, cloze in both languages, German exemplar and factual
recall, translation in both directions), plus the same 12 controls.
Every decomposed surface is validated single-token under the Qwen
tokenizer, cognate pairs are excluded at normalized edit distance
\(\le 2\) after diacritic folding (30 of 132 candidates excluded, with the
exclusion report released alongside the bank), and calls on the surface
language use a lexical criterion, since the two languages share a script.
Jacobian lenses fitted on English and on
German (identical recipe, seeds, and 100-prompt fitting budget) diverge
in their top-1 token readings on 73/90 prompts while the same SRP
feature stays dominant on 84/90 (0.93, CI [0.86, 0.97]), against floors of
1/204 (unrelated token) and 0/270 (shuffled pairing).
Per family the counts are German antonym 7/9, German cloze 14/14, English cloze
15/15, exemplar 7/10, factual recall 8/8, translation DE\(\to\)EN
17/17 and EN\(\to\)DE 16/17, with controls at 5/6 (digits) and 6/6 (proper
nouns). Within a single lens, one feature carries both surface forms on
62/90 (English lens) and 53/90 (German lens) prompts, weaker than in the
English--Chinese study and reported as measured.

\paragraph{One prompt read by three fitted lenses.}
Because the two banks share a model and a basis, one hidden state \(h_\ell\) can be
read by three transports at once, which the pairwise design does not
display. We take \texttt{antonym\_de\_01}, the German analogue of the
published antonym case (``Das Gegenteil von hoch ist
tief. Das Gegenteil von klein ist'', model answer
\texttt{groß}), and read its frozen states with the Jacobian lenses fitted on
English, Chinese, and German
(\Cref{tab:app-cross-lens-de-antonym-layers}). At layer 21 all three
return formatting tokens, and at layer 24 each fills its ranking with its own
fitting language. The English lens reads \texttt{large}, \texttt{
large}, \texttt{-large}, \texttt{Large}, the German lens reads
\texttt{groß}, \texttt{große}, \texttt{large}, \texttt{großes}, and the Chinese
lens reads \texttt{large} followed by
\begin{CJK*}{UTF8}{gbsn}大了, 巨大, 大的\end{CJK*}. The same shift is
visible on a single token, since the transported logit for \texttt{groß}
is 15.16 under the English lens and 15.60 under the Chinese lens against
22.19 under the German lens. By layer 29 the lenses fitted on Chinese
and on German both read \texttt{groß} and the lens fitted on English alone still
reads \texttt{large}.

The decomposition separates the two language pairs here. \texttt{groß} is
carried by f6764 and \texttt{large} by f12474, while \texttt{big} and
\begin{CJK*}{UTF8}{gbsn}大\end{CJK*} share f112. The English and Chinese
realizations of this concept share a readout feature and the English and
German ones do not, which is consistent with the weaker
carrying rate within a single lens reported for German above. That contrast comes from the
row codes of the LM head alone, and this single prompt is reported as a worked
example.

\begin{table}[!tbp]
\caption{\textbf{A German antonym prompt read by three fitted lenses.}
Top-1 token and its transported logit on identical Qwen3.5-9B hidden
states \(h_\ell\) for \texttt{antonym\_de\_01}, whose answer is \texttt{groß}. The
three lenses differ only in fitting corpus.}
\label{tab:app-cross-lens-de-antonym-layers}
\centering
\footnotesize
\setlength{\tabcolsep}{5pt}
\renewcommand{\arraystretch}{1.08}
\begin{tabular}{@{}lccc@{}}
\toprule
Layer & EN-fitted & ZH-fitted & DE-fitted \\
\midrule
24 & \texttt{large} 22.4 & \texttt{large} 21.5 & \texttt{groß} 22.2 \\
26 & \texttt{large} 28.8 & \texttt{large} 29.9 & \texttt{large} 29.1 \\
29 & \texttt{large} 19.0 & \texttt{groß} 19.5 & \texttt{groß} 21.2 \\
\bottomrule
\end{tabular}
\end{table}

\paragraph{A second lens family.}
To test whether corpus conditionality is specific to the Jacobian
construction, we fit a second family of ridge translators
\citep{hoerl1970ridge}, one per layer and without a bias term, into the final-layer
residual basis, in the tuned-lens style \citep{belrose2023tuned} but
solved in closed form in state space, on exactly the same 100-prompt
English and Chinese corpora and layer set as the Jacobian lenses.
Controlling the fitting corpus is what the test requires, so a released
tuned lens is unusable here. It arrives fitted on a corpus we did not
choose, and matching our corpora would mean retraining it in any case. A
closed-form translator keeps that retraining cheap, and because it shares
no fitting procedure with the Jacobian lens, agreement between the two
families cannot be an artifact of a common construction.
Ridge strength is selected per layer on a held-out 10-prompt tail, and
holdout \(R^2\) rises monotonically with depth (0.46--0.92 English,
0.46--0.86 Chinese), and the resulting decoded states match the model's top-1
token on 61--75\% of positions at the two deepest fitted layers. Under
this family the corpus conditionality pattern reproduces (e.g.\
EN\(\to\)ZH translation prompts read Latin under the translator fitted on
English and CJK under the translator fitted on Chinese on 9/10), while the same SRP
feature stays dominant on 67/80 prompts (floors 3/159 and 1/240).
Finally, holding the fitting corpus fixed (English) and varying only
the construction, the Jacobian and ridge lenses read the same dominant
feature on 76/80 prompts with both floors empty (0/159, 0/240), so the
decomposition is invariant to the choice of transport across the two families.

\paragraph{A larger fitting corpus.}
The 100-prompt fitting budget is small, so a divergence between two
lenses could in principle reflect estimation noise. We refit the English, Chinese, and German
Jacobian lenses on 300 prompts each, drawn from the same seeded pools
under an identical recipe, and re-read the same frozen states. The
decomposition does not move, since English--Chinese agreement is 77/80
again and the agreement decision is identical prompt for prompt, with the
same 77 prompts passing, the same three failing, and the counts unchanged
in every family and in both control groups. English--German rises from 84/90 to
85/90. The null floors are unchanged in both pairs (0/159 and 0/240 for
English--Chinese, 1/204 and 0/270 for English--German). The surface reading
does move, since script divergence on the
English--Chinese bank falls from 39/80 to 33/80 and the factual recall
figure quoted in \S\ref{sec:beyond-token} falls from 9/10 to 8/10, while
token divergence on the English--German bank rises from 73/90 to 76/90.
Eleven English--Chinese prompts change their reported surface language
across the two fitting scales, and none changes its dominant feature. If
the divergence were estimation noise that shrinks as the fitting corpus
grows, it would shrink for both pairs, and it does not. What the fitting
budget changes is how often the instrument's surface report flips, while
the readout feature carrying the reading stays fixed.
\Cref{tab:app-cross-lens-extension} collects the agreement rates and both
null floors for the main study and all three extensions.

\begin{table}[!tbp]
\caption{\textbf{Cross-lens extension matrix.} Agreement on the dominant
feature (at least two of layers 21, 24, 26, and 29) with 95\% binomial CIs and the
two null floors (unrelated token / shuffled pairing). Row 1
is the main study, rows 2--3 vary one axis each, row 4 holds the corpus
fixed and varies only the lens construction, and rows 5--6 hold corpus
language and construction fixed while tripling the fitting budget.}
\label{tab:app-cross-lens-extension}
\centering
\scriptsize
\setlength{\tabcolsep}{4.0pt}
\renewcommand{\arraystretch}{1.08}
\begin{tabularx}{\linewidth}{@{}>{\RaggedRight\arraybackslash}Xccc@{}}
\toprule
Comparison & Agreement & 95\% CI & Floors \\
\midrule
EN--ZH, Jacobian & 77/80 & [0.90, 0.99] & 0/159 \(\cdot\) 0/240 \\
EN--DE, Jacobian & 84/90 & [0.86, 0.97] & 1/204 \(\cdot\) 0/270 \\
EN--ZH, ridge translator & 67/80 & [0.74, 0.90] & 3/159 \(\cdot\) 1/240 \\
Jacobian vs ridge, EN & 76/80 & [0.88, 0.98] & 0/159 \(\cdot\) 0/240 \\
EN--ZH, Jacobian, \(n{=}300\) & 77/80 & [0.90, 0.99] & 0/159 \(\cdot\) 0/240 \\
EN--DE, Jacobian, \(n{=}300\) & 85/90 & [0.88, 0.98] & 1/204 \(\cdot\) 0/270 \\
\bottomrule
\end{tabularx}
\end{table}

\paragraph{Caveats and scope.}
The Chinese and German C4 fitting corpora contain incidental English
text, which biases the lens pairs toward each other and makes the
token-level divergence conservative. The study covers one model, now
across two language pairs, two constructions of fitted lenses, and two
fitting scales, and the ridge translator is our minimal instantiation in
the tuned-lens style, simpler than the exact construction of
\citet{belrose2023tuned}. The released records for each case
include every prompt, layer, dominant feature id, and top-1 token reading.
